% Paper 1, §11 — Multimodal extensions
% Page budget: ~5 pages. Anchors: kb/notes/architecture/multimodal-llm-extensions.md;
% kb/excerpts/{alayrac2022-flamingo,llava2023}.md.

\section{Multimodal extensions}
\label{sec:multimodal}
\label{sec:multimodal-extensions}

The architectures of \S\S\ref{sec:kv}--\ref{sec:multitoken} are \glsterm{decoder-only}{decoder-only}
language models over a single token stream. Frontier 2026 systems also
take images, video, and audio. Three integration patterns dominate the
record: gated \glsterm{cross-attention}{cross-attention} insertion into a frozen \glsterm{llm}{LLM}
\citep{alayrac2022-flamingo}, linear projection of vision features into
the \glsterm{llm}{LLM} input-embedding space \citep{llava2023}, and native multimodal
pretraining where vision and text patches share the \glsterm{residual-stream}{residual stream}
from the first token of pretraining
\citep{meta-llama4,internvl3-2025,gemini2-5}. The thesis of this
section is that the projector and the choice of pattern are
genuinely architectural, while the vision encoder and the
freeze-versus-tune schedule are training-pipeline choices that we
defer to Paper~2.

\subsection{The three-piece pattern}
\label{sec:multimodal-three-piece}

Across all three integration patterns the same three pieces appear: a
\emph{vision encoder} $E_V$ that turns an image into patch features, a
\emph{projector} $\pi$ that maps those features into the \glsterm{llm}{LLM}'s residual
width, and the \emph{\glsterm{llm}{LLM} body} that consumes them. The encoder is
typically a CLIP-style or SigLIP-style ViT producing a sequence of
$M$ patch tokens of dimension $d_V$; for a $336{\times}336$ image with
$14{\times}14$ patches, $M{=}576$ and $d_V{=}1024$ in the canonical
LLaVA setup\footnote{KB excerpt: \texttt{kb/excerpts/llava2023.md\#sec-4-1}}
\citep{llava2023}. We write the encoder output and the projection
generically as
\begin{align}
v &= E_V(I) \;\in\; \R^{n_v \times d_V}, \label{eq:multimodal-encoder}\\
\mathbf{X}_{\text{vis}} &= \pi(v) \;\in\; \R^{n_v \times \mathsf{D}}, \label{eq:multimodal-projector}
\end{align}
with $n_v$ the number of visual tokens and $\mathsf{D}$ the \glsterm{llm}{LLM} model
dimension. The interesting structure is in $\pi$ and in how the body
consumes $\mathbf{X}_{\text{vis}}$.

\begin{intuition}
$\mathbf{X}_{\text{vis}}$ is a \emph{soft prompt}: a contiguous block of
$n_v$ vectors in $\R^{\mathsf{D}}$ that the \glsterm{llm}{LLM} treats as if they were
output of its input embedding. The vision encoder is a tokenizer for a
non-text modality, $\pi$ is the embedding lookup, and the \glsterm{llm}{LLM} body is
unchanged from \S\ref{sec:transformer-overview}. Returning to math:
the residual-stream initial state becomes $\mathbf{X}^0 =
[E_{\text{in}}(\mathbf{t}_{1:k}); \pi(E_V(I)); E_{\text{in}}(\mathbf{t}_{k+1:N})]$
with the \glsterm{llm}{LLM}'s ordinary update rule applied above it.
\end{intuition}

This soft-prompt framing is exactly the LLaVA pattern. The other two
patterns deviate by either sending $\mathbf{X}_{\text{vis}}$ through new
\emph{\glsterm{cross-attention}{cross-attention}} layers rather than through the \glsterm{residual-stream}{residual stream}
(Flamingo) or by training $E_V$ and the \glsterm{llm}{LLM} body jointly from scratch on
multimodal data (native multimodal). The pieces are the same; the
plumbing differs.

\subsection{Cross-attention bridge: Flamingo}
\label{sec:multimodal-flamingo}

\citet{alayrac2022-flamingo} introduced the first widely-adopted
multimodal extension by inserting new trainable blocks between the
existing layers of a frozen \glsterm{llm}{LLM}. Two pieces are added: a
\emph{Perceiver Resampler} that compresses the variable-length encoder
output $v$ into a fixed-length set of $k$ latent tokens
$\mathbf{V} \in \R^{k \times \mathsf{D}}$ (typically $k{=}64$),
and \emph{gated \glsterm{cross-attention}{cross-attention} dense} blocks
inserted between every few frozen \glsterm{llm}{LLM} blocks%
\footnote{KB excerpt: \texttt{kb/excerpts/alayrac2022-flamingo.md\#sec-2-1}}.
The inserted block has the form
\begin{align}
\mathbf{y} &= \mathbf{x} + \tanh(\alpha_{\text{xattn}}) \cdot \mathrm{XAttn}(\mathbf{x},\mathbf{V}), \label{eq:flamingo-xattn}\\
\mathbf{z} &= \mathbf{y} + \tanh(\alpha_{\text{ffw}}) \cdot \mathrm{FFW}(\mathbf{y}),  \label{eq:flamingo-ffw}
\end{align}
where $\mathbf{x}$ is the language \glsterm{hidden-state}{hidden state} at the insertion point,
and $\alpha_{\text{xattn}}, \alpha_{\text{ffw}}$ are learnable scalars
\emph{initialized to zero}\footnote{KB excerpt:
\texttt{kb/excerpts/alayrac2022-flamingo.md\#sec-2-1-gated}}
\citep{alayrac2022-flamingo}. With $\tanh(0){=}0$ the inserted block
is the identity at initialization: the conditioned model exactly
reproduces the frozen base LM, and as training proceeds the gates open
and visual information enters the language stream.

The Perceiver Resampler is itself a small \glsterm{transformer}{Transformer} with $k$ learned
\glsterm{query}{query} latents that cross-attend to the encoder output $v$, producing a
fixed-size compressed visual representation regardless of input image
or video length. Crucially, both the vision encoder and the base LM
remain frozen for the duration of training; only the Resampler and the
inserted gated \glsterm{cross-attention}{cross-attention} blocks update
\citep[\S 2.1]{alayrac2022-flamingo}.

\subsection{Linear projection: LLaVA}
\label{sec:multimodal-llava}

\citet{llava2023} took the opposite approach: change nothing about the
\glsterm{llm}{LLM}'s architecture, and replace Flamingo's Resampler-plus-gated-XAttn
machinery with a single trainable linear matrix. With $Z_v = g(X_v)$
the CLIP ViT-L/14 grid features for image $X_v$, LLaVA defines
\begin{equation}
H_v = W \cdot Z_v, \qquad W \in \R^{\mathsf{D} \times d_V}, \label{eq:llava}
\end{equation}
and concatenates the resulting visual tokens $H_v$ with the textual
input as ordinary embedding-space vectors%
\footnote{KB excerpt: \texttt{kb/excerpts/llava2023.md\#sec-4-1}, \glsterm{data-mixing-law}{Eq. 1}}.
This is the entire architectural intervention. The \glsterm{llm}{LLM} body sees a
sequence that begins with $n_v$ visual ``tokens'' and continues with
text; its forward pass is identical to a unimodal \glsterm{llm}{LLM}'s.

Training proceeds in two stages\footnote{KB excerpt:
\texttt{kb/excerpts/llava2023.md\#sec-4-2}}. \emph{Stage 1: \glsterm{feature}{feature}
alignment}, where both the CLIP encoder $g$ and the \glsterm{llm}{LLM} body $f_\phi$
are frozen and only $W$ trains on $595\text{K}$ filtered CC3M image-
text pairs --- the projection learns to produce visual tokens that
look, to the frozen \glsterm{llm}{LLM}, like word embeddings. \emph{Stage 2:
end-to-end fine-tuning}, where the encoder stays frozen but
$\theta = \{W, \phi\}$ trains together on $158\text{K}$ GPT-4-generated
visual instruction examples. Later LLaVA-1.5 replaced the linear $W$
with a 2-layer MLP for small additional gains; the architectural
premise is unchanged.

\begin{analogy}
Flamingo treats the \glsterm{llm}{LLM} as sealed and bolts new vision-conditioning
hardware onto its sides. LLaVA treats the \glsterm{llm}{LLM}'s input-embedding space
as a public API: anything that can be projected into $\R^{\mathsf{D}}$
becomes a token. Returning to math: Flamingo introduces a new operator
$\mathrm{XAttn}(\cdot, \mathbf{V})$ that mixes vision into the residual
stream via Eqs.~\eqref{eq:flamingo-xattn}--\eqref{eq:flamingo-ffw};
LLaVA only changes the input sequence to the unmodified \glsterm{llm}{LLM} via
Eq.~\eqref{eq:llava}.
\end{analogy}

LLaVA's parametric simplicity is the reason it dominated the open
multimodal stack from 2024 onward: a single linear matrix is fast to
train, the \glsterm{llm}{LLM} body needs no surgery, and ablations on data and
encoder choice (CLIP vs.\ SigLIP, image resolution, dynamic patching)
become cheap. Llama 3.2-Vision, InternVL, and the Qwen2-VL line all
build on this LLaVA recipe before diverging on resolution policy and
positional encoding.

\subsection{Native multimodal: Gemini, GPT-4o, Claude, Llama 4, InternVL3}
\label{sec:multimodal-native}

The 2025--2026 frontier moves past adapter patterns. The pretraining
corpus itself is multimodal --- interleaved web pages with text,
images, video frames, and audio --- and the model is trained from
scratch on it \citep{meta-llama4,internvl3-2025,gemini2-5}.
Architecturally this is a standard \glsterm{decoder-only}{decoder-only} \glsterm{transformer}{Transformer} where the
\emph{input embedding is multi-source}: text tokens flow through
$E_{\text{in}}^{\text{text}}$, image patches through a small
ViT-like $E_{\text{in}}^{\text{vis}}$, and audio frames through
$E_{\text{in}}^{\text{audio}}$, all producing $\mathsf{D}$-dimensional
vectors that share the \glsterm{residual-stream}{residual stream}. The body is unchanged; the
``\glsterm{vocabulary}{vocabulary}'' has been generalized from text tokens alone to a union
of token sources. \deepencite{gpt-4o-card §architecture}
\deepencite{claude-3-5-sonnet §architecture}

InternVL3 reports that joint pretraining outperforms an adapter
baseline at matched compute on internal benchmarks
\citep{internvl3-2025}; Llama 4 makes a similar claim
\citep{meta-llama4}. The case for native multimodal at the frontier is
empirical and lab-internal; the architectural commitment is that no
explicit cross-modality bridge ($\mathrm{XAttn}$, Q-Former, or
projector) is needed once the \glsterm{residual-stream}{residual stream} has seen both modalities
since pretraining step zero.

\begin{forumsignal}
Vendor reports (Gemini 2.5, Llama 4, InternVL3) consistently claim
native multimodal pretraining produces better quality at smaller active
parameter count than adapter-based alternatives at the same compute.
The published evidence is internal benchmark comparisons against each
lab's own adapter baselines; no controlled cross-lab study with data,
compute, and post-training fixed has appeared. Open replication at
frontier scale is absent --- treat as forum signal pending
reproduction~\citep{internvl3-2025,meta-llama4}.
\end{forumsignal}

\subsection{Where multimodality is an architecture question}
\label{sec:multimodal-arch-vs-not}

The framing of this paper is that the modern \glsterm{transformer}{Transformer} is a small
set of architectural choices; multimodality adds two genuinely new
ones and cleanly externalizes a third.

\begin{enumerate}
\item \textbf{The integration pattern.} \glsterm{cross-attention}{Cross-attention} bridge
(Flamingo, Eqs.~\eqref{eq:flamingo-xattn}--\eqref{eq:flamingo-ffw}),
prefix-token projection (LLaVA, Eq.~\eqref{eq:llava}), or unified-token
native pretraining are three architecturally distinct choices. They
differ in which weights see vision-conditioning gradients and at what
depth in the network vision can interact with text. This is the
load-bearing axis.

\item \textbf{The projector design.} Whether $\pi$ is a single linear
matrix, a 2-layer MLP, a Q-Former, or a Perceiver Resampler determines
the visual token count $n_v$, the parameter cost, and how aggressively
visual information is compressed before it reaches the residual
stream. Within the prefix-token pattern this is the dominant variable.

\item \textbf{The vision encoder family.} CLIP ViT-L/14, SigLIP,
SigLIP2, Qwen-VL ViT, InternVL ViT --- the choice of $E_V$ is
architectural \emph{at the boundary} between the multimodal interface
and the unimodal vision literature. Inside the \glsterm{llm}{LLM} body its
contribution is summarized by $d_V$ and $n_v$ at the encoder output;
the rest is a separate field's design space.
\end{enumerate}

What is \emph{not} a §\ref{sec:multimodal} architecture question:

\begin{itemize}
\item \emph{Frozen-versus-tuned \glsterm{llm}{LLM} body.} LLaVA's two-stage recipe
freezes then unfreezes the \glsterm{llm}{LLM} body; Flamingo keeps it frozen
throughout; native multimodal trains everything jointly from scratch.
This is a training-pipeline choice with the same architecture under
all three regimes; we defer it to Paper~2 (cf.\ Paper~2's SFT and
multimodal-data-mixture sections).
\item \emph{Visual instruction-tuning data.} LLaVA's
$158\text{K}$ GPT-4-generated visual conversations
\citep[\S 3]{llava2023} are a synthetic-data engineering choice, not
an architectural one --- again Paper~2.
\item \emph{Multi-axis positional encoding.} \glsterm{m-rope}{M-RoPE} and \glsterm{interleaved-mrope}{Interleaved-MRoPE}
extend \glsterm{rope}{RoPE} to 2D and 3D coordinates for image and video patches; the
mechanism is positional and inherits the analysis of
\S\ref{sec:position}, so we treat it there rather than here.
\end{itemize}

The cleanest architectural summary is: vision-language extension
adds one new axis (integration pattern), refines one existing
component (the input-embedding side, now multi-source), and otherwise
leaves the choices of \S\S\ref{sec:kv}--\ref{sec:norm} untouched.

\medskip

A multimodal \glsterm{llm}{LLM} is mostly a unimodal \glsterm{llm}{LLM} with a generalized
input-embedding stage, plus a single new architectural axis for how
visual tokens enter the \glsterm{residual-stream}{residual stream}. The remaining components ---
KV-sharing, attention implementation, \glsterm{ffn}{FFN} family, normalization,
position encoding --- are the same choices Paper~1 has been
documenting throughout. \S\ref{sec:inference-adjuncts} now turns to
the inference-time companions of these choices: KV-cache lifecycle,
quantization regime, and structured-output decoding, which constrain
which architectural choices §§\ref{sec:kv}--\ref{sec:multitoken} can
afford in deployment.
